\documentclass[a4paper,12pt]{article}

\usepackage[T2A]{fontenc}
\usepackage[utf8]{inputenc}
\usepackage[russian]{babel}
\usepackage{amsmath,amssymb,mathtools}
\usepackage{helvet}
\renewcommand{\familydefault}{\sfdefault}
\usepackage{icomma}
\usepackage[a4paper,left=20mm,right=20mm,top=20mm,bottom=22mm]{geometry}
\usepackage{microtype}
\usepackage{xcolor}
\usepackage{tikz}
\usetikzlibrary{calc,arrows.meta,positioning,shapes.geometric}
\usepackage{pgfplots}
\pgfplotsset{compat=1.18}
\usepackage{tabularx,booktabs}
\usepackage{enumitem}
\usepackage{hyperref}
\usepackage{parskip}
\usepackage{fancyhdr}
\usepackage{setspace}
\onehalfspacing

\definecolor{accent}{HTML}{167C80}
\definecolor{darkaccent}{HTML}{0D5558}
\definecolor{textgray}{HTML}{2F3437}
\definecolor{softgray}{HTML}{EEF3F3}
\definecolor{warning}{HTML}{7A4E2D}

\hypersetup{
  colorlinks=true,
  linkcolor=darkaccent,
  urlcolor=darkaccent,
  pdftitle={Графики функций: метод узловых точек — чек-лист},
  pdfauthor={Лёвин Артём Александрович}
}

\setlength{\parindent}{0pt}
\setlength{\headheight}{15pt}
\setlist[itemize]{leftmargin=6mm,itemsep=3pt,topsep=3pt}
\pagestyle{fancy}
\fancyhf{}
\fancyhead[L]{\small Графики функций: метод узловых точек}
\fancyhead[R]{\small Григорий Архипов}
\fancyfoot[C]{\small\thepage}
\renewcommand{\headrulewidth}{0pt}

\newcommand{\heading}[1]{\section*{\color{darkaccent}#1}}
\newcommand{\subheading}[1]{\subsection*{\color{darkaccent}#1}}
\newcommand{\checkitem}[1]{%
  \par\noindent
  \makebox[7mm][l]{\raisebox{0.15ex}{$\square$}}%
  \begin{minipage}[t]{\dimexpr\linewidth-7mm\relax}
    #1
  \end{minipage}\par\vspace{2mm}
}
\newcommand{\key}[1]{\textbf{\color{darkaccent}#1}}
\newcommand{\thinrule}{\par\vspace{1mm}{\color{accent!55}\hrule height 0.5pt}\vspace{3mm}}

\tikzset{
  terminal/.style={ellipse,draw=accent,thick,align=center,text width=5.7cm,minimum height=11mm,font=\small},
  action/.style={rectangle,rounded corners=2mm,draw=accent,thick,align=center,text width=11.4cm,minimum height=12mm,font=\small,inner sep=3mm},
  decision/.style={diamond,aspect=3.1,draw=accent,thick,align=center,text width=6.5cm,font=\small,inner sep=1.6mm},
  arrow/.style={-{Stealth[length=3mm]},thick,draw=accent}
}

\begin{document}
\color{textgray}
\raggedright

\hypersetup{pageanchor=false}
\begin{titlepage}
\centering
\vspace*{22mm}
{\large Профильная математика\par}
\vspace{12mm}
{\Huge\bfseries Чек-лист\par}
\vspace{8mm}
{\LARGE Графики функций\par}
\vspace{3mm}
{\Large Метод узловых точек\par}
\vspace{10mm}
{\large Линейная функция, пересечение графиков, парабола и гипербола\par}
\vspace{18mm}
{\large Григорий Архипов\par}
\vspace{5mm}
{\large 7 октября 2026 года\par}
\vfill
{\normalsize Лёвин Артём Александрович\\эксклюзивно для Григория Архипова\par}
\vspace{14mm}
\begin{tikzpicture}
  \draw[accent,line width=0.8pt] (0,0) .. controls (3.0,0.55) and (7.2,-0.55) .. (11.2,0);
\end{tikzpicture}
\vspace*{5mm}
\end{titlepage}
\hypersetup{pageanchor=true}

\heading{1. Главная идея: восстанавливаем формулу по графику}

\checkitem{Я помню: значение на вертикальной оси можно обозначать как $y$ или как $f(x)$. Для точки графика запись $(x_0;y_0)$ означает $y_0=f(x_0)$.}

\checkitem{Я умею находить \key{узловые точки}: это точки графика, у которых обе координаты читаются точно и являются целыми числами.}

\checkitem{Я беру столько узловых точек, сколько неизвестных коэффициентов нужно определить. Если неизвестны $k$ и $b$, обычно достаточно двух точек.}

\checkitem{Если подходящих точек много, я выбираю самые удобные: с координатой $0$ либо с небольшими по модулю координатами. Это сокращает вычисления и уменьшает риск арифметической ошибки.}

\checkitem{Для каждой выбранной точки подставляю её координаты в формулу функции. Получаю уравнение; несколько неизвестных коэффициентов дают систему уравнений.}

\subheading{Пример: линейная функция}
Пусть график функции $f(x)=kx+b$ проходит через точки $(-1;-3)$ и $(3;4)$. Тогда
\[
\begin{cases}
-3=-k+b,\\
4=3k+b.
\end{cases}
\]
Вычитая первое уравнение из второго, получаем $7=4k$, поэтому
\[
k=\frac74,\qquad b=-\frac54.
\]
Следовательно,
\[
f(x)=\frac74x-\frac54,
\qquad
f(-5)=-10.
\]

\thinrule
\textbf{Самопроверка формулы.} Подставьте координаты обеих исходных точек обратно. Если хотя бы одна точка формуле не удовлетворяет, вычисления нужно перепроверить.

\clearpage
\heading{2. Универсальный алгоритм метода узловых точек}

\begin{center}
\begin{tikzpicture}[node distance=10mm]
  \node[terminal] (start) {Начало: дана формула с неизвестными коэффициентами и график};
  \node[action,below=of start] (count) {Определить, какие коэффициенты неизвестны, и посчитать их количество};
  \node[action,below=of count] (points) {Выбрать столько точных узловых точек, сколько неизвестных коэффициентов};
  \node[action,below=of points] (subst) {Подставить координаты каждой точки в исходную формулу};
  \node[action,below=of subst] (solve) {Решить полученную систему и найти коэффициенты};
  \node[action,below=of solve] (formula) {Собрать формулу функции и проверить её по исходным точкам};
  \node[decision,below=12mm of formula] (task) {Что требуется в условии?};
  \node[action,below left=15mm and -8mm of task,text width=5.2cm] (value) {Найти значение $f(x_0)$: подставить заданный $x_0$};
  \node[action,below right=15mm and -8mm of task,text width=5.2cm] (arg) {Найти аргумент: составить уравнение $f(x)=y_0$ и решить его};
  \node[terminal,below=24mm of task] (finish) {Проверить смысл ответа и знак};

  \draw[arrow] (start)--(count);
  \draw[arrow] (count)--(points);
  \draw[arrow] (points)--(subst);
  \draw[arrow] (subst)--(solve);
  \draw[arrow] (solve)--(formula);
  \draw[arrow] (formula)--(task);
  \draw[arrow] (task)--node[above left,font=\scriptsize]{значение}(value);
  \draw[arrow] (task)--node[above right,font=\scriptsize]{аргумент}(arg);
  \draw[arrow] (value.south) |- (finish.west);
  \draw[arrow] (arg.south) |- (finish.east);
\end{tikzpicture}
\end{center}

\clearpage
\heading{3. Линейная функция: $f(x)=kx+b$}

\checkitem{Я узнаю $k$ как \key{угловой коэффициент}. Его знак показывает направление: при $k>0$ прямая возрастает, при $k<0$ убывает.}

\checkitem{Я использую связь
\[
k=\tg\alpha=\frac{\Delta y}{\Delta x},
\]
где $\alpha$ — угол между прямой и горизонтальным направлением. Для вычисления удобно построить прямоугольный треугольник по двум точным точкам прямой.}

\checkitem{Я помню смысл свободного коэффициента:
\[
b=f(0).
\]
То есть $b$ — ордината точки пересечения прямой с осью $Oy$.}

\checkitem{Если точка пересечения с $Oy$ читается неточно, сначала можно найти $k$, затем взять одну узловую точку $(x_0;y_0)$ и вычислить
\[
b=y_0-kx_0.
\]}

\subheading{Визуальный смысл коэффициентов}
\begin{center}
\begin{tikzpicture}
\begin{axis}[
  width=12.5cm,height=7.2cm,
  axis lines=middle,
  xmin=-4,xmax=4,ymin=-5,ymax=7,
  xtick={-3,-2,-1,0,1,2,3},
  ytick={-4,-2,0,2,4,6},
  grid=both,
  grid style={line width=.1pt, draw=gray!25},
  major grid style={line width=.2pt,draw=gray!35},
  xlabel={$x$},ylabel={$y$},
  clip=false,
  tick label style={font=\small}
]
\addplot[domain=-4:4,samples=100,accent,very thick]{1.5*x+1};
\addplot[only marks,mark=*,mark size=2.2pt,accent] coordinates {(0,1) (2,4)};
\node[anchor=west,font=\small] at (axis cs:0.15,1.05) {$b=1$};
\draw[accent,thick] (axis cs:0,1)--(axis cs:2,1)--(axis cs:2,4);
\node[font=\small] at (axis cs:1,0.55) {$\Delta x=2$};
\node[font=\small,anchor=west] at (axis cs:2.05,2.5) {$\Delta y=3$};
\node[font=\small,anchor=west] at (axis cs:1.1,5.7) {$k=\frac{3}{2}$};
\end{axis}
\end{tikzpicture}
\end{center}

\thinrule
\textbf{Контроль здравого смысла.} После вычислений сопоставьте знак $k$ и положение прямой на рисунке, а значение $b$ — с пересечением оси $Oy$.

\clearpage
\heading{4. Пересечение двух графиков}

Пусть заданы две функции
\[
f_1(x)=k_1x+b_1,
\qquad
f_2(x)=k_2x+b_2.
\]

\checkitem{Я сначала восстанавливаю обе формулы: методом узловых точек или быстрыми приёмами для линейной функции.}

\checkitem{Для точки пересечения выполняется
\[
f_1(x)=f_2(x).
\]
Решение этого уравнения даёт \key{абсциссу} точки пересечения.}

\checkitem{Чтобы получить \key{ординату}, найденный $x$ можно подставить в любую из двух функций. овпадают.}

\subheading{Пример}
Пусть
\[
f_1(x)=4x+9,
\qquad
f_2(x)=\frac32x-\frac72.
\]
Тогда
\[
4x+9=\frac32x-\frac72.
\]
Умножим на $2$:
\[
8x+18=3x-7,
\qquad 5x=-25,
\qquad x=-5.
\]
Подстановка в любую формулу даёт
\[
y=-11.
\]
Точка пересечения: $(-5;-11)$.

\thinrule
\textbf{Внимание к формулировке.} «Найти $f(-5)$», «найти $x$, при котором $f(x)=y_0$» и «найти абсциссу точки пересечения» требуют разных последних шагов, хотя восстановление формулы может быть одинаковым.

\clearpage
\heading{5. Парабола: $f(x)=ax^2+bx+c$}

\checkitem{Я использую тот же принцип: число выбранных узловых точек соответствует числу неизвестных коэффициентов. Если один коэффициент уже дан, уравнений потребуется меньше.}

\checkitem{Свободный коэффициент параболы читается так же, как у прямой:
\[
c=f(0).
\]
Если график пересекает ось $Oy$ в точной точке, $c$ находится сразу.}

\checkitem{Для вершины параболы с абсциссой $x_v$ выполняется
\[
x_v=-\frac{b}{2a},
\qquad
b=-2ax_v.
\]
Если $a$ и $x_v$ известны точно, коэффициент $b$ можно получить без системы.}

\checkitem{Если вершина $V(x_v;y_v)$ и точка при $x=x_v\pm1$ читается точно, то
\[
a=f(x_v\pm1)-y_v.
\]
При сдвиге от вершины на одну единицу по $x$ вертикальное изменение равно $a$.}

\checkitem{Если координаты вершины или соседней точки читаются неточно, я использую метод узловых точек.}

\subheading{Пример с двумя неизвестными коэффициентами}
Пусть
\[
f(x)=2x^2+bx+c
\]
и график проходит через точки $(0;-4)$ и $(1;1)$. Тогда
\[
c=-4,
\]
а из второй точки
\[
1=2\cdot1^2+b\cdot1-4,
\qquad b=3.
\]
Итак,
\[
f(x)=2x^2+3x-4.
\]

\textbf{Быстрый пример с вершиной.} Если $a=1$ и $x_v=-1$, то $b=-2\cdot1\cdot(-1)=2$. При $c=-1$ получаем
\[
f(x)=x^2+2x-1.
\]

\clearpage
\heading{6. Гипербола и другие знакомые формы}

\checkitem{Если в условии формула гиперболы уже дана и в ней два неизвестных коэффициента, я выбираю две удобные узловые точки и подставляю их в эту формулу.}

\checkitem{Я не пытаюсь угадывать коэффициенты только по внешнему виду кривой, когда рисунок не даёт точных значений. Коэффициенты должны следовать из координат и формулы.}

\checkitem{После нахождения коэффициентов я обязательно подставляю исходные точки обратно и проверяю совпадение с графиком.}

\checkitem{Универсальная идея остаётся той же для разных типов функций: известная форма функции $+$ точные точки графика $\rightarrow$ уравнения для коэффициентов $\rightarrow$ восстановленная формула.}

\heading{7. Что особенно важно помнить}

\begin{tabularx}{\linewidth}{@{}>{\bfseries\color{darkaccent}}p{0.31\linewidth}X@{}}
\toprule
Ситуация & Действие \\
\midrule
Нужно найти коэффициенты & Выберите столько точных точек, сколько неизвестных коэффициентов, и составьте систему. \\
Прямая $kx+b$ & $k=\Delta y/\Delta x$, $b=f(0)$. \\
Два графика пересекаются & Приравняйте значения функций: $f_1(x)=f_2(x)$. \\
Нужна ордината пересечения & Подставьте найденный $x$ в любую из двух функций. \\
Парабола & Используйте $c=f(0)$ и $x_v=-b/(2a)$, когда данные читаются точно. \\
График убывает & Для прямой угловой коэффициент $k<0$. \\
Получился странный результат & Сопоставьте формулу с рисунком и проверьте исходные точки. \\
\bottomrule
\end{tabularx}

\vspace{7mm}
\thinrule
\textbf{Итоговая схема:}
\[
\begin{aligned}
\text{точные точки}
&\longrightarrow \text{уравнения для коэффициентов}
\longrightarrow \text{формула функции}\\
&\longrightarrow \text{ответ на конкретный вопрос}.
\end{aligned}
\]

\end{document}
